% Discussion
% Interpret your results, discuss implications, acknowledge limitations

\label{sec:discussion}

\subsection{Interpretation of Results}

The findings indicate that Wikipedia’s openness allows for self-correction of bias over time but does not fully prevent power concentration. The results support the hypothesis that diversity of participation correlates with higher neutrality, consistent with deliberative democratic theory~\cite{klemp2010deliberation, shi2019value}. However, dominance by experienced editors on contentious topics suggests that governance structures can both enable and constrain democratic participation.

\subsection{Implications}

This research contributes to understanding how large-scale online collaboration can balance democratic ideals with the need for content reliability. For Wikipedia, these findings underscore the importance of policies and tools that promote inclusive participation. For the broader field of digital democracy, this study illustrates how algorithmic and governance mechanisms interact to shape knowledge production.

\subsection{Limitations}

Limitations include the scope of analyzed topics, the reliance on automated sentiment and toxicity models (which may misclassify nuanced statements), and the focus on English-language Wikipedia only. Future iterations should expand to multilingual datasets and incorporate deeper qualitative analysis of deliberative exchanges.

\subsection{Future Work}

Future work will refine the bias detection model using transformer-based language models (e.g., BERT or RoBERTa) to improve contextual understanding of bias. Additionally, temporal analysis of participation networks could reveal how editor turnover affects consensus stability. Expanding the dataset to include different language Wikipedias could further illuminate cultural variations in governance and neutrality.
